% ============================================================
% Latent Profiles of Academic Achievement Among Rural Students
% in Peru: An Integrated EFA-IRT-LCA Analysis of PISA Data
% Target: WAIMLAp 2026 -- Springer ACSAR series
% Template: LLNCS (llncs.cls v2.21)
% References: BibTeX splncs04
% ============================================================

\documentclass[runningheads]{llncs}

%% ── Packages ─────────────────────────────────────────────────
\usepackage[T1]{fontenc}
\usepackage[utf8]{inputenc}
\usepackage{graphicx}
\usepackage{amsmath}
\usepackage{amssymb}
\usepackage{url}
\usepackage{booktabs}
\usepackage{multirow}
\usepackage{array}
\usepackage{makecell}
\usepackage{svg}
% \usepackage{hyperref}
\usepackage[hidelinks]{hyperref}
\renewcommand\UrlFont{\color{black}\rmfamily}
\emergencystretch=1.5em

%% ── Abbreviated running head ─────────────────────────────────
%% Title is long; provide a short version for running heads
\newcommand{\shorttitle}{Latent Profiles of Academic Achievement in Rural Peru}

%%%%%%%%%%%%%%%%%%%%%%%%%%%%%%%%%%%%%%%%%%%%%%%%%%%%%%%%%%%%%%%%
\begin{document}
%%%%%%%%%%%%%%%%%%%%%%%%%%%%%%%%%%%%%%%%%%%%%%%%%%%%%%%%%%%%%%%%

\title{Latent Profiles of Academic Achievement Among Rural
Students in Peru: An Integrated EFA--IRT--LCA Analysis of
PISA Data}

\titlerunning{\shorttitle}

\author{Kevin Maximiliano Arenas Cancapa\thanks{Corresponding
author: \email{a20245026@pucp.edu.pe}} \and
Eric Raphael Huiza Pereyra}

\authorrunning{K.M. Arenas Cancapa and E.R. Huiza Pereyra}

\institute{Graduate Program in Computer Science,
Pontifical Catholic University of Peru, Lima, Peru\\
\email{a20245026@pucp.edu.pe}}

\maketitle

%% ── Abstract ─────────────────────────────────────────────────
\begin{abstract}
This research has as its purpose to determine and describe the
latent factors related to the academic performance of secondary
school students from rural areas of Peru, based on the microdata
corresponding to PISA 2022. For this, a methodological pipeline
was developed that combined Exploratory Factor Analysis (EFA),
Item Response Theory (IRT) and Latent Class Analysis (LCA). The
procedure was applied both to the national sample, made up of
6,968 students, and to a rural subsample of 3,720 students. In
the rural setting, the EFA identified 35 latent factors, which
were subsequently refined through IRT calibration from 129 items
with a high discrimination capacity. Through the LCA, four
profiles of rural students were identified, which presented
statistically significant differences in their academic
performance (\(F=327.43\), \(p<.001\)). The first profile
(20.6\% of rural students) corresponded to high-performing
students, 93.5\% of whom attended private schools. The other
profiles comprised two groups of intermediate performance,
differentiated by financial capital and by a limited creative
potential (25.2\% and 17.2\%), as well as a majority group of
low performance (37.0\%), characterized by the accumulation of
disadvantages in the family environment. The proposed
methodology was contrasted with PISA 2018 microdata, whose
results supported its applicability across different evaluation
cycles. Together, the results show that rural educational
difficulties vary according to their mechanisms and profiles,
for which reason they require differentiated and contextualized
interventions.

\keywords{PISA 2022 \and Latent Class Analysis \and Item
Response Theory \and Exploratory Factor Analysis \and Rural
Education \and Academic Achievement \and Peru \and Educational
Inequality \and Latent Variables \and Educational Data Science}
\end{abstract}

%% ─────────────────────────────────────────────────────────────
\section{Introduction}
%% ─────────────────────────────────────────────────────────────

Academic performance in secondary education constitutes a highly
important indicator both of educational quality and of the
long-term social development of an individual and of the country
to which they belong. In Peru, the results of the 2022 cycle of
the Programme for International Student Assessment (PISA) ---
carried out by the Organisation for Economic Co-operation and
Development (OECD) --- revealed that approximately 47\% of
students reached Level~2 or higher in Science, 34\% in
Mathematics and 50\% in Reading, this means that more than half
of the students who participated did not manage to reach basic
levels of competence~\cite{oecd2023pisaPeru}.

Rural schools in Peru operate under structural limitations that
generate patterns of low performance that in essence turn out
different from those observed in urban contexts. This
distinction is important from the analytical point of view:
factors such as the shortage of teachers, limitations in digital
infrastructure and high socioeconomic vulnerability --- which
negatively affect rural settings --- may not operate in the same
way in urban schools, which suggests that analyses carried out
exclusively at the national level can hide relevant mechanisms
belonging to a rural context. The national report of PISA 2022
results empirically supports this differentiation, by
documenting a marked stratification: only 11.0\% of rural
students reached basic competence in Mathematics, 21.8\% in
Science and 21.0\% in Reading, versus 38.7\%, more than 50\% and
more than 50\%, respectively, among urban
students~\cite{minedu2024pisareporte}. The present study takes
motivation from these differences and focuses on the rural
subpopulation as its own analytical unit.

Understanding what explains these performance gaps requires
going beyond the observable variables, to go deeper into the
latent constructs that configure the learning conditions.
Context factors such as family cultural capital, pedagogical
quality, school governance and financial education operate as
non-observable dimensions that influence student performance
without being directly measurable by standard evaluations.
Large-scale datasets like PISA offer enriched contextual
information through questionnaires oriented to students, schools
and teachers, but their analytical potential is usually wasted
when studies are limited to regression approaches or focus on
isolated predictors~\cite{rutkowski2010secondary}.

Addressing these non-observable dimensions requires the use of
methodologies designed for the inference of latent variables.
Recent advances in educational data science have demonstrated
the value of integrating psychometric and machine learning
methods to obtain deeper knowledge from large-scale assessment
data. Exploratory Factor Analysis (EFA) allows reducing the
dimensionality of questionnaire data; Item Response Theory (IRT)
offers an item selection criterion based on discrimination and
difficulty parameters; and Latent Class Analysis (LCA) allows
the identification of distinct subpopulations with shared
contextual profiles. Applied sequentially, these methods
constitute a structured process to move from raw survey data to
interpretable student profiles. However, few studies have
applied this type of integrated processes on PISA data in Latin
American contexts, and the absence is even greater in rural
subpopulations. As far as is known, no previous study was
identified that applies an integrated EFA-IRT-LCA process like
the one presented in this research, focused on rural students in
Peru.

To address this gap, the present research implements a
sequential framework that combines EFA, IRT and LCA applied to
PISA 2022 microdata for Peru, giving priority to the rural
student population. Specifically, this research puts forward
three main objectives: (i)~identify the underlying dimensions in
student, family and institution context variables through EFA;
(ii)~carry out an optimization of the instrument by selecting
items according to IRT discrimination parameters; and
(iii)~outline latent profiles specific to rural students,
evaluating their relationship with performance in each
competency evaluated by PISA: mathematics, reading and science.
Additionally, the differences are compared systematically
between the findings identified for the rural subsample and
those of the complete national sample, which allows
distinguishing the factors that are specific to the rural
context from the structural patterns present throughout the
Peruvian educational system.

In order to evaluate the adaptability of this proposed
analytical framework, the same pipeline was also applied to
PISA 2018 microdata for Peru (\(n=6{,}086\)), obtaining
consistent and reliable results across different evaluation
cycles, without requiring structural adjustments to the
procedure or to the parameters used.

This article has the following structure.
Section~\ref{sec:related} situates the context of the study
within the current academic debate. Section~\ref{sec:method}
details the data sources, the three-stage pipeline and the
estimation procedures. Section~\ref{sec:results} presents the
empirical findings. Section~\ref{sec:discussion} discusses the
implications of these results for educational policy.
Section~\ref{sec:conclusions} closes with the conclusions, the
study's limitations and concrete directions for future research.

%% ─────────────────────────────────────────────────────────────
\section{Related Work}
\label{sec:related}
%% ─────────────────────────────────────────────────────────────

\subsection{Large-Scale Assessments and Educational Inequality
in Latin America}

Educational inequalities are a highly relevant concern globally,
and large-scale evaluations like PISA function as highly useful
instruments to diagnose these disparities and inform
evidence-based policies among participating
countries~\cite{rutkowski2010secondary}. In Latin America and
the Caribbean, PISA results consistently show performance below
the OECD countries' average, a situation that worsened due to
the COVID-19 pandemic. Estimates for the region indicate that
school closures during the pandemic lasted, on average, close to
nine months --- a disruption long enough to be equivalent to the
loss of more than a year of instruction for the average
student~\cite{delprato2025resilience}. Given that this type of
disruption is rarely distributed evenly within an educational
system, students who already faced structural disadvantages ---
this also includes rural areas --- this group can is identified
as those who lost the most learning ground, a pattern that
motivates the subpopulation-level analysis developed in this
present study.

In the particular case of Peru, the urban-rural performance gap
constitutes one of the most structurally entrenched dimensions
of educational inequality. Using data from the Student Census
Evaluation (ECE) 2018, a study carried out with Peruvian
secondary school students found a gap of 63.28 points in
mathematics performance in favor of urban students, of which
82.58\% was explained by observable differences --- mainly
socioeconomic factors and school resources --- between both
groups~\cite{rojasapaza2024urbanrural}. A second layer of
division adds to these geographic disparities: given that
families with economic capacity tend to opt for private
education, the resulting comparison between public and private
students is never entirely clean --- families who withdraw from
the public system are not necessarily equal to those who remain,
in characteristics that the data fail to fully capture. This
selection effect also does not operate uniformly across the
whole country, since the relative performance gap between school
types varies depending on whether the institution is located in
an urban or rural setting~\cite{pagliarani2025schools}, a
pattern that this study takes up by examining school management
type as a discriminant variable of the latent classes within the
rural subsample.

The recovery of learning loss associated with the pandemic has
also been different between school types: improvements in
mathematics derived from post-pandemic educational interventions
in the region were more pronounced in private institutions than
in public ones, and rural schools that experienced longer
closures showed a different recovery trajectory compared to
their urban counterparts~\cite{delprato2025covid}. This
asymmetry raises the question of whether it becomes even more
accentuated when the analysis is centered on a rural
subpopulation, one of the structural dimensions that this study
examines through the identification of latent performance
profiles. Policies aimed directly at the shortage of rural
teachers have also not had more positive results. Monetary
incentives designed to keep teachers in rural areas do improve
retention in the short term, but their subsequent effect on
students' actual learning begins to decrease once the
institutional support that accompanies them is
withdrawn~\cite{castro2022bonuses} --- an early indication that
low rural performance is hardly resolved through
single-dimension interventions, this corresponds to one more
reason to examine the multiple latent factors that interact in
this study's model.

\subsection{Latent Variable Methods in Educational Research}

Exploratory Factor Analysis (EFA) is a tool that is still valid
and fundamental for reducing the dimensionality of large sets of
items and discovering latent dimensions in questionnaire
data~\cite{fabrigar1999efa}. Item Response Theory (IRT) offers a
framework that complements the estimation of discrimination and
difficulty parameters at the item level, this allows for an item
selection based on psychometric
criteria~\cite{vanderlinden2016irt}. Instead of assigning each
student a single ability score, PISA reports competence as a set
of plausible values: several values from a distribution that
reflect how much uncertainty persists about a student's
authentic level of ability, due to limited exposure to the test
items. Analysts working with PISA data are expected to use these
plausible values as the appropriate unit of analysis and not
just a single point estimate~\cite{vondavier2009plausible}, a
convention on which this study draws to construct the composite
performance score used to validate the latent class solution
(the implications of averaging the plausible values are
discussed in Section~\ref{sec:limitations}).

Latent Class Analysis (LCA) extends the latent variable
framework toward a person-centered classification, identifying
subpopulations of individuals who share similar response
patterns~\cite{collins2010latent}. The present study implements
LCA using StepMix~\cite{morin2025stepmix}, an open-source Python
package that implements pseudo-likelihood estimation for
generalized mixture models. These methods are well established
and frequently applied to large-scale assessment data, but
generally in isolation and for different purposes. IRT is at the
core of PISA's own scaling procedures, where it is used to
estimate proficiency and to build questionnaire indices from
predefined item sets~\cite{oecd2024technical}, whereas
person-centred studies usually derive profiles from variables or
indices selected a priori. The contribution of the present
pipeline lies in their sequential integration, in which each
stage prepares the input of the next: EFA identifies latent
dimensions empirically from the full set of questionnaire items;
IRT acts as a psychometric filter that discards items with low
discrimination, which would add noise rather than information to
the classification; and LCA classifies students on this screened
item set. The EFA factors are then reused to interpret the
resulting classes, while the plausible values, which do not
enter the model, serve as an external validation of the
profiles. Moreover, applications of this kind to Latin American
PISA data remain scarce~\cite{delprato2025resilience}, and none,
to our knowledge, has focused on rural subpopulations.

\subsection{Student Profile Analysis and Machine Learning in
Large-Scale Assessments}

Person-centered approaches for the analysis of PISA data
currently gained considerable momentum~\cite{repo2025wellbeing}.
Explainable cluster analysis applied to PISA 2022 data also
showed that student profiles derived from context variables are
capable of achieving high interpretability when combined with
variable importance methods~\cite{alvarezgarcia2024profiles}.

A large-scale transnational analysis of PISA 2022 data, which
worked with more than seventy educational systems, found that
household cultural resources --- commonly measured through book
possession --- consistently manage to predict cognitive
outcomes, although the strength of the relationship is not
uniform: it weakens or strengthens depending on whether the
country represents the OECD or not~\cite{djekourmane2025capital}.
This inconsistency between national contexts constitutes in
itself an argument to consider with caution when applying
findings about cultural capital derived from aggregated
comparisons between countries to a single non-OECD country, such
as Peru --- and even more so to a rural subpopulation within it,
whose socioeconomic conditions may be very different from the
national average used in such studies. This interpretation is
consistent with Bourdieu's foundational
framework~\cite{bourdieu1986capital}, which defines cultural
capital as a set of knowledge, skills and dispositions
transmitted through family socialization, which translate into
educational advantage.

Studies that combine statistical and machine learning methods
demonstrated their usefulness for identifying predictors of
student performance in PISA
datasets~\cite{soares2024predictors};
however, research that was applied making use of interpretable
methods to educational systems in the global south remains
scarce, and most applications concentrate on contexts of more
developed countries, whose realities differ markedly from those
of developing nations~\cite{delprato2025resilience}.

Together, the reviewed literature evidences three analytical
gaps: (1)~the absence of latent profile analysis focused
specifically on rural student populations in Peru; (2)~the
limited application of integrated EFA-IRT-LCA pipelines to PISA
questionnaire data; and (3)~the limited exploration of how
structural factors, such as school management type and creative
potential development, can influence performance profiles in
rural educational contexts.

%% ─────────────────────────────────────────────────────────────
\section{Methodology}
\label{sec:method}
%% ─────────────────────────────────────────────────────────────

\subsection{Data Source and Study Population}

This study uses microdata from the PISA 2022
cycle~\cite{oecd2024technical}. For Peru, the evaluation was
carried out with 6968 students from schools across the
territory. Specifically for this research, the following six
datasets were used: student cognitive data
(\texttt{STU\_COG}), student context questionnaire
(\texttt{STU\_QQQ}), student time data (\texttt{STU\_TIM}),
school context questionnaire (\texttt{SCH\_QQQ}), teacher
context questionnaire (\texttt{TCH\_QQQ}) and classroom or
collective cognitive data (\texttt{CRT\_COG}). Originally, each
dataset was in \texttt{.SAV} format, for this reason it was
necessary to convert it to \texttt{.CSV} format for its handling
and treatment. To interpret the data collection, the PISA
reference document (Codebook) was used. As a next step, the
records were linked through the student (\texttt{CNTSTUID}) and
school (\texttt{CNTSCHID}) identification variables. This
resulted in a combined dataset of 6968 observations and 7147
variables before preprocessing.

Because PISA 2022 lacks official subnational reports for Peru,
the rural subsample was carried out using the school-level
variable \texttt{SC001Q01TA}, which classifies the type of
community to which the school belongs. Two categories were
retained: village, hamlet or rural area (fewer than 3000
inhabitants) and small town (from 3000 to approximately 15,000
inhabitants). This filtering yielded 3720 rural students, which
corresponded to 53.4\% of the national sample. It should be
noted that the rural subsample (\(n = 3720\)) is nested within
the national sample (\(n = 6968\)); the subsequent evaluation
presents these groups as two separate analytical perspectives on
a shared dataset, rather than treating them as autonomous sample
populations. The complete national sample was maintained as a
comparative baseline.

\subsection{Analytical Pipeline Overview}

The workflow consists of four stages: (1)~data preprocessing and
variable reduction; (2)~exploratory factor analysis (EFA) for
the identification of latent dimensions; (3)~item response
theory (IRT) parameter estimation for item selection; and
(4)~latent class analysis (LCA) for student profile
identification. Each step must be integrated with the next
through a systematic reduction of variables. This process was
carried out independently on both samples.

The analysis focuses on PISA 2022 because it is the most recent
cycle available and includes dimensions central to this study,
such as the module on school closures during COVID-19,
questionnaire items related to creative thinking and an expanded
financial literacy questionnaire. As a first application of the
proposed pipeline to Peruvian PISA data, the study prioritises
depth of analysis within a single cycle over breadth across
cycles. PISA 2018, the most recent previous cycle, was used to
assess whether the pipeline could be replicated on a different
dataset without methodological modifications, rather than to
pool data across cycles. Earlier cycles were not included
because changes in questionnaire content and assessment mode
limit the comparability of the item-level inputs on which the
pipeline operates.

To evaluate the generalization of this workflow across PISA
cycles, it was applied to PISA 2018 microdata (\(n = 6086\))
without methodological modifications. It was confirmed that the
rural subsample (\texttt{SC001Q01TA}) was available in PISA 2018
with category definitions identical to those of PISA 2022.
Table~\ref{tab:replication} presents a comparative summary of
the process results in both cycles at the national level.

\begin{table}[!htbp]
\centering
\caption{Pipeline Replication: Comparative Outputs Across PISA
Cycles (National Sample)}
\label{tab:replication}
\begin{tabular}{lcc}
\toprule
\textbf{Metric} &
  \makecell{\textbf{PISA 2022} \\ \textbf{(n=6,968)}} &
  \makecell{\textbf{PISA 2018} \\ \textbf{(n=6,086)}} \\
\midrule
Variables after preprocessing  & 592    & 438    \\
Items after EFA                & 235    & 229    \\
Items after IRT ($a \geq 0.50$)& 152    & 161    \\
Mean IRT discrimination        & 1.029  & 0.954  \\
Optimal LCA classes            & 7      & 10     \\
BIC inflection point           & 8 classes & 11 classes \\
Entropy (optimal solution)     & 0.9796 & 0.9730 \\
ANOVA $F$ (class validation)   & 397.59*** & 681.79*** \\
\bottomrule
\multicolumn{3}{l}{\footnotesize *** \textit{p} $< .001$}
\end{tabular}
\end{table}

To replicate the findings of this research --- both in the
national and rural sample --- the results are available through
a read REST API, composed of 16 endpoints segmented into: EFA
(loaded factor structures), IRT (item calibration parameters),
LCA (class profiles, model selection indices and discriminant
variable classifications) and contextual characterization
(school type and community type distributions by latent class).
All these endpoints support query parameters for both rural and
national samples. The URL of the API can be found in the README
of the repository, and the code, preprocessing scripts and model
specifications required to reproduce the analysis are available
at \url{https://github.com/kevin-arenas/pisa-project} (see Data
and Code Availability).

\begin{figure}[!htbp]
\centering
\includesvg[width=0.6\textwidth]{graphics/fig1_pipeline_v4}
\caption{Analytical process from raw PISA microdata to latent
student profiles. The numbers indicate the number of variables
at each reduction stage for the national and rural samples.}
\label{fig:pipeline}
\end{figure}

\subsection{Data Preprocessing}

The preprocessing stage reduced the initial 7,147-variable
dataset through five sequential filters applied identically to
both samples.

\textbf{Step 1 --- Missing data filtering.} Filtering of missing
data. The elimination of variables with more than 40\% null
values was carried out. According to PISA, this lack of data
occurs in a planned way due to the distribution of questions in
the booklets, this is the reason why the imputation of these
variables is not an appropriate
alternative~\cite{oecd2024technical}.

\textbf{Step 2 --- Design variable exclusion.} The identifiers,
sampling weights and composite indices derived from the OECD
were eliminated, with respect to the plausible-values-type
variables they were separated to use later as external
performance anchors in the LCA validation
step~\cite{vondavier2009plausible}.

\textbf{Step 3 --- Free-text exclusion.} The columns with free
text that had more than 20\% unique values were eliminated for
not being useful for factor analysis.

\textbf{Step 4 --- Continuous variable separation.} The
numeric-type columns with more than 10 unique values were
discarded from the ordinal item set, the reason was that they
were not applicable for the methodologies used later ---
exploratory factor analysis (EFA) and item response theory
(IRT) --- these require ordinal or categorical type inputs.

\textbf{Step 5 --- Near-zero variance elimination.} Finally, the
variables in which a single category covered more than 95\% of
the responses were eliminated.

Table~\ref{tab:preprocessing} presents the reduction summarized
in each stage described in previous paragraphs. For the
variables retained after preprocessing, missing data averaged
7.12\% in the rural sample; median imputation was applied for
exploratory factor analysis (EFA), then mode imputation for item
response theory (IRT) and latent class analysis.

\begin{table}[!htbp]
\centering
\caption{Variable Reduction at Each Preprocessing Stage}
\label{tab:preprocessing}
\begin{tabular}{lcc}
\toprule
\textbf{Stage} &
  \makecell{\textbf{National} \\ \textbf{(n=6,968)}} &
  \makecell{\textbf{Rural} \\ \textbf{(n=3,720)}} \\
\midrule
Original variables         & 7,147 & 7,147 \\
After missing $>$ 40\%     & 1,134 & 1,098 \\
After design var.\ removal &   893 &   857 \\
After continuous sep.      &   648 &   625 \\
After near-zero variance   &   592 &   573 \\
\bottomrule
\end{tabular}
\end{table}

\subsection{Stage 1 --- Exploratory Factor Analysis}

Exploratory factor analysis (EFA) was carried out using the
Python package \texttt{factor\_analyzer}, implementing principal
axis factor analysis with oblimin rotation. Oblique rotation was
preferred because educational constructs --- family support,
school climate, student motivation --- are theoretically and
empirically correlated~\cite{fabrigar1999efa}. Given the high
number of retained variables, the EFA was carried out through
thematic blocks defined by the origin of the questionnaire and
the prefix of the variables, following the practices established
in large-scale research~\cite{rutkowski2010secondary}. The number
of defined blocks was six: student items (\texttt{ST}), three
school subpopulations --- due to the high number of variables
for this prefix (\texttt{SC\_a}, \texttt{SC\_b},
\texttt{SC\_c}) ---, items related to financial literacy
(\texttt{FL}) and a block that grouped the residual variables.
In addition, a block intended for parent-related items retained
only one variable after preprocessing and was excluded, since
factor analysis requires multiple items; this accounts for the
difference between the 573 preprocessed variables and the 572
items analysed in Table~\ref{tab:efa_blocks}. Before extraction,
factorability was evaluated through KMO and Bartlett's test;
blocks with KMO~$< 0.50$ were discarded. Factors were retained
using Kaiser's criterion (eigenvalue $> 1$) and the retained
items had to have a factor loading $\geq 0.40$.

Table~\ref{tab:efa_blocks} shows the previously mentioned blocks
with the results for the rural sample. As can be observed in
Table~\ref{tab:efa_blocks}, School Block A was excluded
(KMO~$= 0.485$) in the rural sample --- unlike the national
sample, where it was processed successfully (KMO~$= 0.809$) ---.
This finding is highly relevant; the governance and evaluation
variables showed an insufficient factor structure in the rural
sample, which probably reflects the homogeneity of rural
institutions within the country.

\begin{table}[!htbp]
\centering
\caption{EFA Results by Questionnaire Block (Rural Sample)}
\label{tab:efa_blocks}
\begin{tabular}{lccccc}
\toprule
\textbf{Block} & \textbf{Items} & \textbf{KMO} &
  \textbf{Status} & \textbf{Factors} &
  \makecell{\textbf{Items retained}} \\
\midrule
Student (ST)     & 136 & 0.809 & Adequate   & 10 & 47 \\
School A (SC\_a) &  99 & 0.485 & Excluded   & -- &  0 \\
School B (SC\_b) &  92 & 0.546 & Acceptable &  6 & 34 \\
School C (SC\_c) & 109 & 0.557 & Acceptable &  6 & 49 \\
Financial (FL)   &  93 & 0.895 & Adequate   &  8 & 51 \\
Other            &  43 & 0.664 & Adequate   &  5 & 17 \\
Parents          &   1 & --    & Excluded (single item) & -- & 0 \\
\midrule
\textbf{Total}   & \textbf{573} & & & \textbf{35} & \textbf{198} \\
\bottomrule
\end{tabular}
\end{table}

Table~\ref{tab:efa_factors} presents the ten student-block
factors with their interpretive labels and primary items.

\begin{table}[!htbp]
\centering
\caption{Student-Block EFA Factors --- Rural Sample
(10 factors, 47 items)}
\label{tab:efa_factors}
\begin{tabular}{clp{3.6cm}}
\toprule
\textbf{F} & \textbf{Interpretive label} & \textbf{Primary items} \\
\midrule
1  & Household socioeconomic resources             & ST254, ST251, ST250 \\
2  & Academic motivation \& self-concept           & ST268Q07--09 \\
3  & Household resources (country-specific items)  & ST250D, ST251D \\
4  & School safety and belonging                   & ST265, ST034 \\
5  & Victimization and school insecurity           & ST038, ST266 \\
6  & Support for learning mathematics              & ST297 \\
7  & Participation in creative school activities   & ST337 \\
8  & Parental educational level                    & ST006, ST008 \\
9  & Beliefs about academic ability                & ST263Q04, ST263Q06 \\
10 & Teacher support for learning                  & ST270 \\
\bottomrule
\end{tabular}
\end{table}

The main structural differences between the national and rural
factor solutions are presented in Table~\ref{tab:efa_comparison},
where the reordering of the main student factor in the context
of each sample is observed: participation in extracurricular
activities (national), household socioeconomic resources
(rural), which evidences the difference between the rural and
national context. In the rural context, household socioeconomic
resources represent the central point of differentiation between
students.

\begin{table}[!htbp]
\centering
\caption{EFA Structural Comparison: National vs.\ Rural Sample}
\label{tab:efa_comparison}
\begin{tabular}{p{2.5cm}p{3.8cm}p{3.8cm}}
\toprule
\textbf{Dimension} & \textbf{National} & \textbf{Rural} \\
\midrule
F1 (student) &
  Extracurricular participation &
  Household socioeconomic resources \\
School Block A &
  Processed (KMO=0.809) &
  Excluded (KMO=0.485) \\
Unique factors &
  Art/creativity attitudes; math self-perception and attitudes &
  Participation in creative school activities; parental
  education; teacher support \\
Items retained & 235 & 198 \\
\bottomrule
\end{tabular}
\end{table}

\subsection{Stage 2 --- Item Response Theory Calibration}

IRT calibration was applied to the 198 items retained from the
EFA using the Python package GIRTH~\cite{eribean2021girth},
which provides marginal maximum likelihood estimation for
unidimensional IRT models. Two models were implemented:

\begin{itemize}
\item \textbf{The Two-Parameter Logistic (2PL):} model for
dichotomous items, estimating discrimination (\(a\)) and
difficulty (\(b\)) parameters.
\item \textbf{Graded Response Model (GRM):} for polytomous items
(3--10 ordered categories); extends 2PL to multiple thresholds.
\end{itemize}

Items with an \(a\) value $\geq 0.50$ were retained, a more
conservative criterion than the commonly cited
\(a \geq 0.35\)~\cite{vanderlinden2016irt}, adopted to ensure
that only highly discriminatory items entered the LCA.
Table~\ref{tab:irt_results} presents the calibration results.
The mean discrimination of the items retained in the rural
sample (1.109) exceeded that of the national sample (1.029),
which confirms that the rural item set captures more marked
individual differences.

\begin{table}[!htbp]
\centering
\caption{IRT Calibration Outcomes (Rural Sample)}
\label{tab:irt_results}
\begin{tabular}{lcccc}
\toprule
\textbf{Model} &
  \makecell{\textbf{Items} \\ \textbf{input}} &
  \makecell{\textbf{Retained}\\\textbf{($a \geq 0.50$)}} &
  \textbf{Discarded} &
  \makecell{\textbf{Mean}~$a$\\\textbf{(retained)}} \\
\midrule
2PL   &  38 &  21 &  17 & 1.878 \\
GRM   & 160 & 108 &  52 & 0.960 \\
\midrule
\textbf{Total} & \textbf{198} & \textbf{129} &
  \textbf{69} & \textbf{1.109} \\
\bottomrule
\end{tabular}
\end{table}

\subsection{Stage 3 --- Latent Class Analysis}

LCA was carried out on the 129 retained items using
StepMix~\cite{morin2025stepmix} with a categorical measurement
model and one-step estimation (\texttt{measurement=}
\texttt{"categorical"}, \texttt{n\_steps=1},
\texttt{max\_iter=500}, \texttt{random\_state=42}). Models with
2 to 8 classes were estimated and compared using four joint
criteria: (1)~Bayesian Information Criterion (BIC; lower =
better fit); (2)~relative entropy ($\geq 0.80$ indicates
adequate classification certainty); (3)~minimum class size
($\geq 5\%$ of the sample, that is, $\geq 186$ students for
rural areas); and (4)~substantive interpretability. Class
profiles were characterized through a two-layer strategy:
Layer~1 used standardized z-score factor scores to describe each
class along the EFA dimensions; Layer~2 validated the classes
against academic performance using the average of all 30
plausible values in mathematics, reading and science as a
composite achievement score. Differences between classes were
evaluated through a one-way analysis of variance (ANOVA).

%% ─────────────────────────────────────────────────────────────
\section{Results}
\label{sec:results}
%% ─────────────────────────────────────────────────────────────

\subsection{Preprocessing and Pipeline Replication}

The process was also applied to PISA 2018 microdata
(\(n = 6086\)) exclusively at the national level, following the
same procedure as the one carried out with PISA 2022 microdata.
Since the purpose of this replication was to verify the
reproducibility of the pipeline rather than to characterise the
rural population in 2018, the rural replication and the
interpretation of the 2018 profiles are left for future work.
Table~\ref{tab:replication} confirms that the EFA-IRT-LCA
sequence is capable of producing coherent and methodologically
consistent results across cycles. The differences between the
retained variables (592 versus 438) reflect an expected and
genuine structural variation: PISA 2022 introduced questionnaire
modules on COVID-19 and expanded the financial literacy items
absent in 2018, which rules out instability of the process. The
consistency of the entropy values (0.9796 versus 0.9730) and the
statistical significance of the differences between classes
(\(F = 397.59\) and \(F = 681.79\), both \(p < 0.001\)) confirm
robust latent class solutions independently of the input cycle.

\subsection{EFA: Latent Factor Structures}

Exploratory factor analysis (EFA) identified 35 latent factors
from 198 retained items in five processable blocks
(Table~\ref{tab:efa_blocks}). Of these 35 identified factors,
four factors from the ``Others'' block (F2-F5) were discarded
due to their administrative nature for the PISA test ---
adaptive module assignment, database version control --- rather
than student or context constructs. It should be noted that F2
also included the school type variable (\texttt{PRIVATESCH})
together with the rest of the administrative indicators, which
may suggest a possible association between school management
type and the assignment of adaptive test modules in the rural
sample: a pattern that requires its own research, but which is
outside the scope of the present study. This exclusion was
interpretive rather than procedural: the items of these four
factors, including \texttt{PRIVATESCH}, remained in the
pipeline, were submitted to the IRT calibration and, when
retained, entered the LCA. They were only omitted from the
substantive characterisation of the classes
(Table~\ref{tab:lca_zscores} and Fig.~\ref{fig:heatmap_rural}),
which is therefore based on the remaining 31 factors.

The student block grouped ten factors
(Table~\ref{tab:efa_factors}), where F1 indicated household
socioeconomic resources as the key factor for the
differentiation of rural students with respect to the national
sample, where the key factor corresponded to the frequency of
extracurricular participation (Table~\ref{tab:efa_comparison}).

\subsection{IRT: Item Discrimination and Selection}

Of the 198 items retained by exploratory factor analysis (EFA),
the estimation through the item response theory (IRT) model
preserved 129 items for latent class analysis (LCA) (retention
rate of 65.2\%). Table~\ref{tab:irt_results} presents the
estimation results and Fig.~\ref{fig:irt_discrimination} shows a
comparison between the national and rural level discrimination
parameters.

\begin{figure}[!htbp]
\centering
\includegraphics[width=\textwidth]{graphics/fig2_irt_disc_v4.png}
\caption{IRT discrimination parameter distribution --- national
vs rural sample}
\label{fig:irt_discrimination}
\end{figure}

\subsection{LCA: Model Selection}

Table~\ref{tab:lca_fit} presents the model fit indices for 2 to
8 classes in the rural sample. The 4-class solution was selected
as the only model that simultaneously satisfied the four
criteria. Models with five or more classes produced at least one
class below the minimum 5\% threshold (\(n < 186\)), which
demonstrates that the sample is not large enough for specific
subgroups within the rural population. In the national sample,
the BIC increased at 8 classes, which confirms that the 7-class
model is more adequate, since each one has more than 5\% of the
cases. (ANOVA \(F = 397.59\), \(p < 0.001\)).

\begin{table}[!htbp]
\centering
\caption{LCA Model Fit --- Rural Sample (2--8 Classes)}
\label{tab:lca_fit}
\begin{tabular}{lcccc}
\toprule
\textbf{Classes} & \textbf{BIC} & \textbf{Entropy} &
  \makecell{\textbf{Min.\ class size (\%)}} &
  \textbf{Viable} \\
\midrule
2 & 844,801 & 0.9997 & --     & $\checkmark$ \\
3 & 831,527 & 0.9640 & --     & $\checkmark$ \\
4 & 819,868 & 0.9753 & 17.2\% & $\checkmark$ \\
5 & 815,989 & 0.9783 &  4.3\% & $\times$ \\
6 & 813,541 & 0.9773 &  4.8\% & $\times$ \\
7 & 811,454 & 0.9821 &  4.0\% & $\times$ \\
8 & 809,780 & 0.9819 &  3.8\% & $\times$ \\
\bottomrule
\multicolumn{5}{l}{\footnotesize
  $\times$ = minimum class size $< 5\%$ threshold ($n < 186$).}
\end{tabular}
\end{table}

\subsection{LCA: Rural Student Profiles}

The plausible values of mathematics, reading and science ---
previously separated in the preprocessing stage --- were used to
identify the average academic performance of each latent class,
which constitutes an external validation, since these values did
not enter the model. Table~\ref{tab:lca_profiles} also reports
the size of each class and the management type of the
institutions (private or public); unlike performance, school
type (\texttt{PRIVATESCH}) was one of the LCA indicators and is
therefore reported descriptively.

\begin{table}[!htbp]
\centering
\caption{Rural Latent Class Profiles: Performance and Context}
\label{tab:lca_profiles}
\begin{tabular}{lcccc}
\toprule
\textbf{Class} &
  \makecell{\textbf{Mean PV}\\\textbf{(pts)}} &
  \textbf{SD} & \textbf{n (\%)} &
  \makecell{\textbf{Private school \%}} \\
\midrule
1 High-performing  & 446.2 & 70.3 & 768 (20.6\%)   & 93.5\% \\
2 Medium           & 398.3 & 67.7 & 938 (25.2\%)   &  0.0\% \\
3 Medium-low       & 382.2 & 70.3 & 639 (17.2\%)   &  0.0\% \\
4 Low-performing   & 354.7 & 59.0 & 1,375 (37.0\%) &  0.0\% \\
\bottomrule
\multicolumn{5}{p{0.8\textwidth}}{\footnotesize ANOVA:
  $F = 327.43$, $p < .001$. PV = mean of 30 plausible values.
  School type (\texttt{PRIVATESCH}) was one of the 129 LCA
  indicators; its distribution across classes partly reflects
  its role in class formation.}
\end{tabular}
\end{table}

The results obtained from latent factors and plausible values
allowed an interpretation of the class profiles at two levels.
At Level~1 (internal characterization), classes are
distinguished by their standardized factor score profiles
(Table~\ref{tab:lca_zscores}); at Level~2, classes are validated
through their association with academic performance --- a
variable that did not participate in the formation of the latent
classes --- using the average of the 30 plausible values in
mathematics, reading and science.

\begin{table}[!htbp]
\centering
\caption{Standardised Factor Scores (z) by Latent Class ---
Dimensions selected for their relevance}
\label{tab:lca_zscores}
\begin{tabular}{lcccc}
\toprule
\textbf{Latent dimension} &
  \textbf{C1} & \textbf{C2} & \textbf{C3} & \textbf{C4} \\
\midrule
School management and governance                 & $+1.76$ & $-0.48$ & $-0.58$ & $-0.39$ \\
Students' creativity and creative potential      & $-0.40$ & $+0.34$ & $-1.12$ & $+0.52$ \\
Financial and digital literacy                   & $+0.38$ & $+0.65$ & $+0.01$ & $-0.66$ \\
Professional development of mathematics teachers & $+0.93$ & $-0.37$ & $+0.18$ & $-0.35$ \\
Household socioeconomic resources                & $+0.78$ & $+0.23$ & $-0.26$ & $-0.47$ \\
Educational conditions during COVID-19 closures  & $+0.80$ & $-0.21$ & $-0.05$ & $-0.28$ \\
Parents' educational attainment                  & $+0.57$ & $+0.03$ & $-0.16$ & $-0.26$ \\
\bottomrule
\multicolumn{5}{l}{\footnotesize
  C1=High-performing; C2=Medium; C3=Medium-low; C4=Low-performing.}
\end{tabular}
\end{table}

While Table~\ref{tab:lca_zscores} presents the seven latent
dimensions with the greatest interpretive relevance for the
classes, the complete characterization of each class is carried
out based on the 31 latent factors, with the purpose of
considering the more subtle patterns specific to each class that
may be relevant to each interpretation. The four profiles are
interpreted as follows:

\textbf{Class~1 --- High-Management Schools with Elevated Family
Capital (768 students, 20.6\%):} This first class is composed of
factors that can be grouped into three contexts: institutional,
family and student. In the institutional context is found the
dimension with the greatest magnitude of the entire table ---
school management and governance (\(z = +1.76\)) ---, accompanied
by a notable investment in teacher professional development in
mathematics (\(z = +0.93\)) and favorable educational conditions
during the COVID-19 closures (\(z = +0.80\)). This combination
suggests that solid institutional management, together with
sustained investment in teacher training, may have facilitated a
more effective response to an adverse context such as the
pandemic.

It is important to note that this class is composed in 93.5\% by
privately managed institutions, unlike the other three classes,
composed entirely of public management. This composition implies
that the institutional values described cannot be attributed
with certainty exclusively to a specific latent profile, since
they are partially confounded with the general effect of
belonging to the private sector. For this reason, the more
reliable comparisons in this study are those made between
Classes 2, 3 and 4, where school management type remains
constant.

Moreover, since \texttt{PRIVATESCH} was among the LCA
indicators, this separation is partly built into the class
solution: what the model reveals is not that school type defines
the class by itself, but that it co-occurs with the strongest
institutional and family advantages of the sample.

From the family context, household socioeconomic level
(\(z = +0.78\)) and parental education level (\(z = +0.57\))
turned out to be determinant in the composition of this class,
consistent with its status as the class with the highest
academic performance according to the plausible values used as
an external validation variable.

It is particularly relevant that, despite being the class with
the highest academic performance, academic motivation and
self-concept (\(z = -0.29\)) and student creativity and creative
potential (\(z = -0.40\)) are located below the general average.
This suggests that the high performance of this class is not
explained by greater academic motivation or by a higher creative
potential attributed to students by their schools, but rather
responds mainly to an accumulated structural and institutional
advantage.

\textbf{Class~2 --- Medium Performers with Financial Literacy as
Protective Capital (938 students, 25.2\%):} This group contains
the traits most related among classes oriented to financial
capital, both at the student level and in the family context ---
financial and digital literacy (\(z = +0.65\)), use of digital
financial services (\(z = +0.53\)) and financial socialization
in the family (\(z = +0.28\)) are located consistently above the
average. Unlike Class~1, the school context does not stand out
positively in school management and governance (\(z = -0.48\))
nor in teacher professional development in mathematics
(\(z = -0.37\)), which could partly explain the low result in
educational conditions during the COVID-19 closures
(\(z = -0.21\)), despite having an institutional preparation for
remote education above average (\(z = +0.15\)). However,
specific support for mathematics learning (\(z = +0.29\)) and
curricular planning in this area (\(z = +0.17\)) remain slightly
above average, suggesting that the institutional weaknesses are
not generalized but concentrated in management and in the
response to the pandemic.

From the student context, this class is the only one of the four
with a level of victimization and school insecurity above
average (\(z = +0.17\)), a distinctive trait that clearly
differentiates it from the others. Student creativity and
creative potential (\(z = +0.34\)) is also located above
average.

Unlike Class~1, this class belongs entirely to the public
sector, which allows more reliable institutional comparisons
against Classes 3 and 4. The general pattern suggests that the
strength of family and student financial capital, together with
household socioeconomic resources slightly above average
(\(z = +0.23\)), co-occur with medium academic performance
(\(M = 398\) pts) despite weak institutional management, a
pattern consistent with a compensatory role that cannot be
tested with cross-sectional data.

\textbf{Class~3 --- Creative Paradox: Low Student Potential
despite Favorable School Climate (639 students, 17.2\%):} This
class is characterized by presenting consistent weaknesses in
the institutional context: school management and governance
below average (\(z = -0.58\)), curricular planning in
mathematics below average (\(z = -0.40\)), and reduced
institutional policies for the use of digital technologies
(\(z = -0.32\)). Curiously, institutional preparation for remote
education (\(z = -0.35\)) --- a value similar in magnitude to
the one observed in Class~1 --- did not translate into
unfavorable educational conditions during the COVID-19 closures,
which remained practically neutral (\(z = -0.05\)). This
relationship turns out striking when compared to Class~2, where
the inverse pattern occurs: institutional preparation above
average did not prevent notably inferior educational response
conditions during the pandemic. This suggests that prior
institutional preparation and the educational response during
the closures are not consistently associated across classes.

In the family context, this class presents household
socioeconomic resources below average (\(z = -0.26\)), together
with a parental education level slightly below average
(\(z = -0.16\)), which together suggests a reduced level of
family capital.

In the student context, the most marked paradox of the four
classes is identified: despite having the only positive value
among all classes in the ``school creativity climate'' dimension
(\(z = +0.55\)), the students of this class reflect the most
negative value of the entire sample in ``student creativity and
creative potential'' (\(z = -1.12\)). Since both dimensions
derive from the school questionnaire, this discrepancy could
suggest a gap, within the principals' own reports, between the
creative climate attributed to the institution and the creative
potential observed in its students; however, this interpretation
should be taken as a hypothesis to explore in future research,
given that the available data do not allow establishing a causal
relationship between both dimensions.

\textbf{Class~4 --- High Motivation and Creativity despite Family
and Financial Capital Disadvantage (1375 students, 37.0\%):} This
group contains the lowest performance level among the four
classes and the most numerous population of students. In the
institutional context, school management and governance
(\(z = -0.39\)) and teacher professional development in
mathematics (\(z = -0.35\)) are located below average. However,
unlike Class~3, the impact of teacher evaluation (\(z = +0.16\))
remains slightly above average, as does curricular planning in
mathematics (\(z = +0.19\)), institutional policies for the use
of digital technologies (\(z = +0.12\)) and, more notably,
institutional preparation for remote education (\(z = +0.29\))
--- the highest positive institutional value among all classes
for this latent dimension. However, educational response
conditions during the COVID-19 closures are located below
average (\(z = -0.28\)).

In the family context, parental education level (\(z = -0.26\))
and household socioeconomic resources (\(z = -0.47\)) are below
average, with more negative values than those observed in
Class~3, which evidences the differentiated weight of this
context among classes.

From the student context, the most relevant finding of this
class is identified: student creativity and creative potential
(\(z = +0.52\)) and academic motivation and self-concept
(\(z = +0.16\)) are not only above average, but represent the
highest values recorded for both dimensions among the four
classes. This suggests that, while creativity and motivation are
usually associated with better academic performance, in this
rural population these characteristics do not manage to
translate into favorable performance when the family and
institutional contexts present significant disadvantages ---
unlike what was observed in Classes 2 and 3.

Finally, this class presents the lowest values in the financial
block among all the evaluated dimensions --- both in financial
and digital literacy (\(z = -0.66\)) and in family financial
socialization (\(z = -0.30\)) and exposure to financial
education at school (\(z = -0.18\)) ---. It is especially
notable that, together with household socioeconomic resources,
the financial dimensions are among the dimensions that most
consistently follow the same order as academic performance among
Classes 2, 3 and 4: the higher the performance, the higher the
financial score. This suggests that the base of financial
knowledge and practices could have a particularly close
relationship with academic performance in the rural context.

Fig.~\ref{fig:context_class_rural} synthesizes this contextual
characterization, presenting the complete distribution of school
management type and community type for the four latent classes,
and allows visualizing jointly the institutional and geographic
composition patterns described previously.

\begin{figure}[!htbp]
\centering
\includegraphics[width=\textwidth]{graphics/fig3_context_classes_rural_v4.png}
\caption{Contextual characterisation of latent classes --- rural
sample (\(n=3{,}720\)). School type (left) was one of the LCA
indicators; community type (right) did not enter the model.
Legend labels: Privada = private; P\'ublica = public; Pueblo
peque\~no = small town.}
\label{fig:context_class_rural}
\end{figure}

\subsection{National vs.\ Rural Profile Comparison}

Both solutions --- the national one with 7 classes and the rural
one with 4 --- have similarities with respect to the high levels
in their latent dimensions represented in each block, such as
household socioeconomic resources --- \(z=+0.76\) in the
national sample and \(z=+0.78\) in the rural sample --- in the
student block, as well as governance and curricular management
--- ``Institutional quality management'' \(z=+0.67\) in national
and ``School management and governance'' \(z=+1.76\) in rural
---. In addition to presenting the highest levels with respect to
the mean, these values also correspond in both cases to the
classes with the highest performance, which suggests that the
characteristics in both samples present a similar general
pattern.

However, when observing the highest scores of each block, a
difference can be noted in the dimensions that have the greatest
impact for each sample. In the national sample, within the
student block, the most representative dimension is
``Participation in extracurricular activities'' (\(z=+1.05\),
Class~3), unlike the rural subsample, where the most
representative latent dimension is ``Household socioeconomic
resources'' (\(z=+0.78\), Class~1).

Generally, in the school blocks of the national sample, high
dimensions are observed regarding governance, autonomy and
school resources that highlight the impact of these
characteristics on performance; the highest dimension in this
block is ``Autonomy and school administrative management''
(\(z=+1.38\)) in Class~2. In the school blocks of the rural
sample, on the other hand, we have ``School management and
governance'' (\(z=+1.76\)) as the most representative dimension,
followed by dimensions related to teacher training and the
school's response conditions during COVID-19. This implies a
difference in the impact of each characteristic for each group
within the school block: in the rural sample, the most relevant
factors were teacher training and educational response to
COVID-19, while in the national sample, in addition to these
characteristics, the school's institutional resources and
conditions were also relevant.

Finally, in the financial block, different patterns are observed
between samples. In the national sample, Class~1 concentrates
the highest values --- ``Financial and digital competence''
(\(z=+0.47\)) and ``Digital financial behavior'' (\(z=+0.48\))
---, with a general decreasing trend toward the lower-performing
classes. In the rural sample, on the other hand, the most
representative levels are not concentrated in the
best-performing class, but in Class~2 --- ``Financial and
digital literacy'' (\(z=+0.65\)) and ``Use of digital financial
services'' (\(z=+0.53\)) ---, which suggests that, unlike the
national pattern, financial capital in the rural context is not
necessarily associated with the highest academic performance
class, but rather operates as a differentiating resource within
the public sector.

\begin{figure}[!htbp]
\centering
\includegraphics[height=0.82\textheight,keepaspectratio]{graphics/fig4_heatmap_rural_v4.png}
\caption{Latent class profiles by EFA-derived latent variables.
Rural sample (\(n=3{,}720\))}
\label{fig:heatmap_rural}
\end{figure}

%% ─────────────────────────────────────────────────────────────
\section{Discussion}
\label{sec:discussion}
%% ─────────────────────────────────────────────────────────────

\subsection{School Management Type and Institutional Advantage in
Rural Peru}

In rural Peru, school management type is closely tied to student
profiles and academic results. Because school type was one of
the LCA indicators, the near-complete separation between Class~1
(93.5\% private) and Classes 2--4 (entirely public) should not
be read as an independent finding, but as evidence that school
type structures the profiles jointly with other advantages. What
the model does reveal is that private schooling co-occurs with
the highest score in the school management and governance factor
(\(z = +1.76\)) --- a magnitude considerably greater than that
observed at the national level (\(z = +1.60\) and \(+1.61\) in
the two highest-performing classes) --- together with high
teacher professional development and family capital, which
affirms and adds to previous evidence on public-private
inequalities in Latin America. In rural contexts, where the
accessibility of students to private institutions is
considerably reduced but not nonexistent, the students who do
access them belong to a small clearly advantaged group, whose
difference in performance compared to their public school peers
reflects accumulated advantages in teaching quality, cultural
capital and institutional resources. This finding has high and
direct implications for educational policy: strategies aimed at
improving rural educational quality must be capable of
distinguishing between substantially different institutional
environments in which rural school students are educated.

This result should also be read alongside the exclusion of
School Block A in the rural sample (KMO~$= 0.485$). School type,
a structural attribute of the institution, must be distinguished
from the internal governance and evaluation practices that Block
A was intended to capture. The lack of factor structure in this
block suggests that rural public schools, which operate under a
common administrative framework, show too little variation in
these practices to form latent dimensions; part of this
dimension is nonetheless represented by the school management
and governance factor that emerged in School Block C.
Consequently, the results indicate which type of institution
rural students attend, but not how internal governance differs
among rural public schools. For the same reason, the most
informative comparisons of this study are those among Classes
2--4, where school type is constant and the differences arise
from other dimensions.

\subsection{Creative Potential as a Differentiating Dimension
Among Public Schools}

The identification of Class~3 as a distinctive profile
characterized by a highly restricted student creative potential
(\(z = -1.12\)) reveals that, within the rural public sector, not
all students face the same conditions. While Classes 2 and 4 show
values close to average or slightly positive in this same
dimension, Class~3 is characterized by the most extreme value of
the entire factorial characterization, this turns out to be
particularly interesting, since the same class presents the only
positive value among the four classes in the ``school creativity
climate'' factor (\(z = +0.55\)) --- a measure more related to
how the institution declares its environment in relation to the
students ---, which suggests a possible gap, within the school's
own reports, between the creative climate attributed to the
institution and the creative potential observed in its students.
This subgroup (17.2\% of rural students) distinguishes itself
from Class~4 precisely because its lower performance is
associated with this specific restriction of creative potential,
rather than with a generalized deprivation of family capital,
this suggests that programs aimed at strengthening institutions
in a targeted way could have differentiated effectiveness for
this group compared to the broader structural interventions that
Class~4 would need.

\subsection{Financial Literacy as Compensatory Capital in Rural
Contexts}

The emergence of financial literacy as a dimension that allows
differentiating between Classes 2 and 4 --- it is important to
mention that both are composed entirely of public school
students --- suggests a mechanism that has received little
attention in the Latin American educational literature. In rural
economic contexts where household financial management
constitutes a daily practical necessity among family members,
financial knowledge could function as a form of practical
capital that can be leveraged to favor cognitive development
without this depending solely on formal school quality. This
interpretation is consistent with Bourdieu's framework of forms
of capital, extended here to include practical financial literacy
as a specific variant of cultural capital in this context. Future
research should examine whether interventions aimed at financial
education generate indirect benefits on academic performance in
rural contexts.

\subsection{Accumulated Disadvantage and the Limits of School
Adaptation}

The paradox identified in Class~4 --- where above-average student
creative potential (\(z = +0.52\)) is present alongside the
lowest academic performance (\(M = 355\) pts) and the highest
deficiencies in family capital --- points to a fundamental
limitation of interventions oriented only at the school setting
when facing contexts of high vulnerability. This class, the most
numerous of the sample (37\% of rural students), shows relatively
positive institutional indicators --- institutional preparation
for remote education, curricular planning and impact of teacher
evaluation slightly above average --- without this being able to
translate into better academic performance. That does not imply
that these institutional factors are not important, but that
their association with performance may depend on a base level of
family resources that a significant proportion of rural Peruvian
students currently does not possess.

\subsection{Methodological Contribution}

Beyond its findings obtained in the research, the study
demonstrates the high feasibility and analytical value of an
integrated EFA-IRT-LCA pipeline applied to PISA microdata. Its
contribution lies less in the individual techniques, which are
well established, than in their sequential integration. Unlike
analyses based on OECD-derived composite indices, which were
deliberately excluded during preprocessing, the pipeline builds
the classification from psychometrically screened item-level
data, allowing context-specific dimensions, such as financial
literacy or creative potential in the rural sample, to emerge
empirically rather than being imposed by predefined constructs.
The sequential logic --- EFA for dimension identification, IRT
for psychometric item selection, LCA for person-centered
classification --- allows the production of interpretable
profiles, grounded both in statistical rigor and theoretical
coherence. The effectiveness of the pipeline (reducing 7,147
initial variables to 129 items of high analytical quality)
addresses a challenge that is still present in the secondary
analysis of large-scale assessments like PISA. The consistent
application of this same pipeline both for the comparison between
national and rural samples and for the validation between the
PISA 2022 and 2018 cycles reinforces its methodological soundness
as an analytical framework that can be replicated.

%% ─────────────────────────────────────────────────────────────
\section{Conclusions, Limitations, and Future Work}
\label{sec:conclusions}
%% ─────────────────────────────────────────────────────────────

\subsection{Conclusions}

This study applied an integrated EFA--IRT--LCA pipeline to PISA
2022 microdata to identify latent academic performance profiles
among rural students in Peru. Four distinct profiles were
identified, ranging from a minority of high-performing students
in private schools (20.6\%) to a numerous majority of vulnerable
students (37\%) highly represented by accumulated disadvantages
in the family, educational and financial domains. School
management type emerged as a structural axis of the profiles: the
private-school minority concentrated in the highest-performing
class together with the strongest institutional and family
advantages, a pattern more pronounced in the rural than in the
national context. Secondary determinants such as: student
creative potential, financial literacy and family cultural
capital; further differentiated public school students into
subgroups with clearly distinct intervention needs.

Among the dimensions examined, school management type, together
with the latent factors of student creative potential, family
cultural capital and financial literacy, emerged as the
dimensions with the greatest interpretive relevance for
differentiating student profiles. These factors are not abstract
statistical constructs: they correspond to identifiable and
measurable conditions in students' school and family
environments.

The four latent profiles crucially reveal that the rural public
school population is not homogeneous in its disadvantage.
Financial literacy appears as a potentially compensatory
characteristic in Class~2, restricted creative potential defines
Class~3, and accumulated family capital deficiencies characterize
Class~4. This internal differentiation within the public sector
is a finding that aggregated statistics would not reveal, and has
direct implications for the design and targeting of educational
interventions. From a computational and data science point of
view, this study demonstrates that an integrated psychometric
pipeline, applied to large-scale assessment data, is capable of
producing empirically grounded profiles that go considerably
beyond what conventional analytical approaches on the same data
can produce.

\subsection{Limitations}
\label{sec:limitations}

Several limitations should be considered when interpreting these
findings. First, PISA 2022 does not offer a subnational breakdown
for Peru; the rural subsample was defined through community size
(\texttt{SC001Q01TA}), a key variable that approximates but does
not precisely capture specific regional characteristics. In
particular, the analysis does not distinguish between natural
regions (coast, Andes and Amazon), nor does it examine linguistic
or cultural factors, such as the use of indigenous languages,
which are central to rural education in Peru. Future research
should replicate this methodology using instruments with regional
coverage, such as Peru's Student Census Evaluation (ECE). Second,
PISA's modular design produces missing data of a structural
nature that cannot be completely resolved through imputation; the
40\% threshold was established based on a methodological
criterion, but its selection remains partially arbitrary. Third,
the 4-class solution, although supported by multiple criteria,
reflects the statistical power limitations of the rural subsample
(\(n = 3{,}720\)); a larger rural sample could support a finer
profile differentiation. Fourth, PISA's cross-sectional design
prevents causal inference; the relationships identified between
the identified latent profiles and academic outcomes should be
understood as descriptive, without assuming causal effects.
Fifth, the School Block A variables could not be included in the
rural analysis due to their insufficient factor structure
(KMO~$= 0.485$), leaving uncharacterized a highly relevant
dimension of school functioning in the rural context, something
that did not happen in the national sample.

Accordingly, the role of school management type reported here
refers to the public--private distinction, which was itself an
LCA indicator, rather than to internal governance practices. A
sensitivity analysis excluding \texttt{PRIVATESCH} from the LCA
would be needed to determine whether the separation of the
private-school class persists independently of this variable.

In addition, the number of factors was determined using Kaiser's
criterion (eigenvalue $> 1$), which is known to overestimate the
number of factors. Although the IRT stage subsequently filtered
out weakly discriminating items, complementary retention
procedures such as parallel analysis or Velicer's MAP test would
provide stronger support for the 35-factor solution.

Finally, class-level performance was summarised using the
average of the 30 plausible values. While this yields unbiased
class means, it underestimates measurement variance; estimating
the ANOVA on each plausible value and combining the results would
provide more accurate inferential statistics, although the
magnitude of the observed differences (\(F = 327.43\)) makes it
unlikely to alter the conclusions. Likewise, the analyses were
unweighted, so the results describe the sample rather than
providing population estimates.

\subsection{Future Work}

The application of the pipeline to the ECE's longitudinal data
could allow a regional breakdown and causal inference through
panel methods. The extension to other LAC region countries with
rural subsamples in PISA 2022 would also allow evaluating the
generalization of the identified profiles. An integration of
geospatial analysis to examine the geographic distribution of the
latent classes within rural Peru could allow identifying priority
intervention areas. The development of a confirmatory version of
the pipeline (CFA--2PL--LCA) would allow testing hypotheses
derived from these exploratory findings in later PISA cycles. The
public API developed for this study currently serves the outputs
of both PISA 2022 and PISA 2018; extending it with the
interpretation of the 2018 profiles would establish a
longitudinal interface for tracking the evolution of latent
profiles across evaluation periods.

Re-estimating the latent classes without school type as an
indicator would clarify whether the private-school profile
emerges from the remaining institutional and family dimensions
alone.

Applying the pipeline to the rural subsample of PISA 2018 and to
other PISA cycles would allow assessing the temporal stability of
the rural profiles. Causal questions, however, would require
longitudinal data such as the ECE, since successive PISA cycles
are independent cross-sectional samples.

Future work will also examine the robustness of the factor
solution using parallel analysis and the MAP test.

\subsection{Data and Code Availability}

The PISA 2022 and 2018 microdata are publicly available from the
OECD~\cite{oecd2024technical}. The code and scripts for each
stage of the pipeline, together with the model specifications,
are available at \url{https://github.com/kevin-arenas/pisa-project}.
The processed outputs of the EFA, IRT and LCA stages for PISA 2022
and PISA 2018 are accessible through a public REST API, whose URL
can be found in the README of the repository.

%% ─────────────────────────────────────────────────────────────
%% Credits (required by Springer)
%% ─────────────────────────────────────────────────────────────
\begin{credits}
\subsubsection{\ackname}
The authors used AI-assisted writing tools to support
manuscript preparation, including grammar correction,
language refinement, style improvement, and editorial
feedback. All AI-generated content was critically reviewed
and revised by the authors prior to inclusion. The authors
take full responsibility for the accuracy and integrity
of the work.

\subsubsection{\discintname}
The authors have no competing interests to declare that are
relevant to the content of this article.

\subsubsection{Author Contributions}
K.M. Arenas Cancapa: Conceptualization, Data curation,
Formal analysis, Investigation, Methodology, Software,
Visualization, Writing --- original draft.
E.R. Huiza Pereyra: Supervision, Writing --- review
and editing.
\end{credits}

%% ─────────────────────────────────────────────────────────────
%% Bibliography
%% ─────────────────────────────────────────────────────────────
\bibliographystyle{splncs04}
\bibliography{references_springer}

\end{document}
